(i) On the sixteen, the even half of Proposition \claimref{standard-model.hyp-s-slot-statistics} (the halving the Pati-Salam stage forces, Theorem \claimref{synthesis.neutral-singlet-bit-decided-by-gauge-algebra-leaving-u5}(iii); the Dirac-neutrino entries need it), with the pairing of Lemma \claimref{standard-model.invariant-pairing-five-modes}: by Theorem \claimref{standard-model.yukawa-pattern-two-rung-module}, $w_{4}$ couples exactly $(L, e^{c})$ and $(Q, d^{c})$ and $v_{4}$ exactly $(L, \nu^{c})$ and $(Q, u^{c})$, each with symmetric form. For a diagonal $g$ with value $g(\varphi)$ on a basis state, $B(\psi, w_{4}g\varphi) = g(\varphi)B(\psi, w_{4}\varphi)$, so the symmetric part on a coupled pair $\{\psi, \varphi\}$ is $\tfrac12(g(\psi) + g(\varphi))$ times the $10$ entry, and likewise for $v_{4}$. By Proposition \claimref{standard-model.generation-ten-under-pati-salam} and the formula $T_{3R} = \tfrac12(N_{R} - 1)$ of Proposition \claimref{standard-model.automorphisms-preserve-slot-pair-split-are}, the doublets $L$, $Q$ have $T_{3R} = 0$, the singlets $e^{c}$, $d^{c}$ have $T_{3R} = \tfrac12$ and $\nu^{c}$, $u^{c}$ have $T_{3R} = -\tfrac12$, while $B - L$ is $1$ on $e^{c}$ and $\nu^{c}$ and $-\tfrac13$ on $d^{c}$ and $u^{c}$. So for $g = h(T_{3R}, B - L)$ the down-type lepton weight is $\tfrac12(h(0, -1) + h(\tfrac12, 1))$ and the up-type one $\tfrac12(h(0, -1) + h(-\tfrac12, 1))$, and on quarks the same with $(\tfrac13, -\tfrac13)$. Writing $h = h_{e} + h_{o}$ with $h_{e}$ even and $h_{o}$ odd in $T_{3R}$, $h_{o}(0, \cdot) = 0$ and $h_{o}(-\tfrac12, \cdot) = -h_{o}(\tfrac12, \cdot)$, which gives $e + o$ and $e - o$. The magnitudes of the $10$ entries agree in the two sectors under the conjugate reading (Corollary \claimref{standard-model.one-vev-vector-relations-breaking-pattern}: both are $|b|$), so equal weights mean equal masses.

(ii) The split. By Proposition \claimref{synthesis.georgi-jarlskog-square-of-majorana-element}(i), $TT^{\dagger} = \prod_{i}(\tfrac12 + \beta_{i})$ with $\beta_{i} = n_{i} - \tfrac12$, and $G_{4} = \tfrac18\sum_{i < l}\beta_{i}\beta_{l}$, each $\beta_{i}\beta_{l}$ a monomial supported on the modes $i$, $l$. The pair-only part is $P_{4} = \tfrac18\sum_{i < l \le 3}\beta_{i}\beta_{l} = \tfrac{1}{16}\bigl((\beta_{1} + \beta_{2} + \beta_{3})^{2} - \tfrac34\bigr) = 4\mu_{P}^{2} - \tfrac{3}{64}$ (using $\beta_{i}^{2} = \tfrac14$), which with $\mu_{P} = -\tfrac{3}{16}(B - L)$ is $\tfrac{3}{64}(3(B - L)^{2} - 1)$, and on pair number $k = 0, 1, 2, 3$ takes the values $\tfrac{3}{32}, -\tfrac{1}{32}, -\tfrac{1}{32}, \tfrac{3}{32}$, equal to $\tfrac{3}{32}(B - L)(-1)^{k}$. The rung-only part is $R_{4} = \tfrac18\beta_{4}\beta_{5} = \tfrac{1}{32}(-1)^{N_{R}} = -\tfrac{1}{32}\chi$, since $(-1)^{n} = -2\beta$ for each mode. The rest is $X_{4}$ (loc. cit. (ii)). The same split is Theorem \claimref{synthesis.georgi-jarlskog-is-rung-vector-times-pair-part-of-majorana-square}(i). Weights: $P_{4}$ is a function of $B - L$, with values $\tfrac{3}{32}$ at $|B - L| = 1$ and $-\tfrac{1}{32}$ at $\tfrac13$, the same on both states of each pair, so its weights are those values in both sectors. $\mu_{R} = \tfrac18(\beta_{4} + \beta_{5}) = \tfrac18(N_{R} - 1) = \tfrac14T_{3R}$ is $0$ on doublets and $\pm\tfrac18$ on the singlets, so by (i) its weights are $\tfrac{1}{16}$ down and $-\tfrac{1}{16}$ up on both sectors. $X_{4} = 8\mu_{R}\mu_{P}$ with $\mu_{P} = -\tfrac{3}{16}(B - L)$ (Proposition \claimref{synthesis.georgi-jarlskog-majorana-moment-map}(i)) is $-\tfrac{3}{8}T_{3R}(B - L)$: zero on doublets, $-\tfrac{3}{16}$ on $e^{c}$, $\tfrac{1}{16}$ on $d^{c}$, $\tfrac{3}{16}$ on $\nu^{c}$ and $-\tfrac{1}{16}$ on $u^{c}$; halving gives the stated weights. $R_{4} = -\tfrac{1}{32}\chi$ takes opposite values on the doublet and the singlet of each pair, and $\mu_{P}$, a multiple of $B - L$, opposite values because every pair is $B - L$ neutral (Proposition \claimref{families.doublet-direction-module-georgi-jarlskog-clebsch}(iii)); both give antisymmetric forms only.

(iii) Adding (ii): down $(\tfrac{3}{32} - \tfrac{3}{32}, -\tfrac{1}{32} + \tfrac{1}{32}) = (0, 0)$, up $(\tfrac{3}{16}, -\tfrac{1}{16})$. The up-type form of Proposition \claimref{families.doublet-direction-module-georgi-jarlskog-clebsch}(iv) has moduli $1$ on $(L, \nu^{c})$ and $\tfrac13$ on $(Q, u^{c})$; on the sixteen $v_{4}P_{4} = -\tfrac{3}{32}v_{4}(B - L)\chi$ by the same parity identity as for $w_{4}$ ($\chi = -(-1)^{N_{P}}$ on the sixteen), and $X_{4}$ has the same up weights as $P_{4}$, so the up entries of $G_{4}$ are $-\tfrac{3}{16}$ times that form. Through degree four, $TT^{\dagger}$ is $\tfrac{1}{32} + \tfrac12(\mu_{P} + \mu_{R}) + G_{4}$ (loc. cit. (i)); by (ii) the down weights are $\tfrac{1}{32} + \tfrac{1}{32} + 0 = \tfrac{1}{16}$ on both sectors and the up weights $(\tfrac{1}{32} - \tfrac{1}{32} + \tfrac{3}{16}, \tfrac{1}{32} - \tfrac{1}{32} - \tfrac{1}{16})$. Finally $TT^{\dagger}$ is the projector onto the full state (loc. cit. (i), with the Majorana element of Proposition \claimref{standard-model.majorana-element}), which is odd, so $TT^{\dagger}\varphi = 0$ for $\varphi$ in the sixteen and the whole square gives no entry.

The check computes all of this on the $32 \times 32$ rational module: the sector weights of $w_{4}g$ and $v_{4}g$ against the $10$ forms for four $T_{3R}$-even and two $T_{3R}$-odd weightings with hand-tabulated values, the parts $P_{4}$, $X_{4}$, $R_{4}$ of the expanded square, $\mu_{R}$, $\mu_{P}$, $G_{4}$, the truncation through degree four, and the whole square.
